\Needspace{6\baselineskip}
\begin{center}\hypertarget{native-249-article177}{}\pdfbookmark[2]{177.}{article-177}177.\end{center}
Daraus folgt unmittelbar, dass das wahrscheinlichste System von Werten der Größen $p$, $q$, $r$, $s$ usw. dasjenige ist, bei dem $\Omega$ den größten Wert annimmt. Es ist daher aus den $\nu$ Gleichungen $\frac{\mathrm{d}\Omega}{\mathrm{d}p}=0$, $\frac{\mathrm{d}\Omega}{\mathrm{d}q}=0$, $\frac{\mathrm{d}\Omega}{\mathrm{d}r}=0$, $\frac{\mathrm{d}\Omega}{\mathrm{d}s}=0$ usw. zu ermitteln. Setzt man $M-V=v$, $M'-V'=v'$, $M''-V''=v''$ usw. sowie $\frac{\mathrm{d}\varphi\Delta}{\varphi\Delta\,.\,\mathrm{d}\Delta}=\varphi'\Delta$, so nehmen diese Gleichungen die folgende Form an:
\[
\begin{gathered}
\frac{\mathrm{d}v}{\mathrm{d}p}\varphi'v+
\frac{\mathrm{d}v'}{\mathrm{d}p}\varphi'v'+
\frac{\mathrm{d}v''}{\mathrm{d}p}\varphi'v''+\text{usw.}=0\\
\frac{\mathrm{d}v}{\mathrm{d}q}\varphi'v+
\frac{\mathrm{d}v'}{\mathrm{d}q}\varphi'v'+
\frac{\mathrm{d}v''}{\mathrm{d}q}\varphi'v''+\text{usw.}=0\\
\frac{\mathrm{d}v}{\mathrm{d}r}\varphi'v+
\frac{\mathrm{d}v'}{\mathrm{d}r}\varphi'v'+
\frac{\mathrm{d}v''}{\mathrm{d}r}\varphi'v''+\text{usw.}=0\\
\frac{\mathrm{d}v}{\mathrm{d}s}\varphi'v+
\frac{\mathrm{d}v'}{\mathrm{d}s}\varphi'v'+
\frac{\mathrm{d}v''}{\mathrm{d}s}\varphi'v''+\text{usw.}=0\\
\text{usw.}
\end{gathered}
\]

Hieraus lässt sich durch Elimination eine vollständig bestimmte Lösung der Aufgabe ableiten, sobald die Beschaffenheit der Funktion $\varphi'$ bekannt ist. Da sie sich nicht a priori bestimmen lässt, gehen wir die Sache von der anderen Seite an und fragen, welche \SourceLeaf{250}Funktion der geläufige Grundsatz, dessen Vorzug allgemein anerkannt wird, eigentlich gleichsam stillschweigend voraussetzt. Man pflegt nämlich die Annahme als Axiom anzusehen, dass bei einer durch mehrere unmittelbare Beobachtungen unter gleichen Umständen und mit gleicher Sorgfalt bestimmten Größe das arithmetische Mittel aller beobachteten Werte den wahrscheinlichsten Wert ergibt, wenn schon nicht mit absoluter Strenge, so doch wenigstens näherungsweise, sodass es stets am sichersten ist, sich daran zu halten. Setzen wir also $V=V'=V''=\text{usw.}=p$, so muss allgemein $\varphi'(M-p)+\varphi'(M'-p)+\varphi'(M''-p)+\text{usw.}=0$ gelten, wenn man für $p$ den Wert $\frac{1}{\mu}(M+M'+M''+\text{usw.})$ einsetzt, für welche positive ganze Zahl $\mu$ auch steht. Nehmen wir also $M'=M''=\text{usw.}=M-\mu N$ an, so gilt allgemein, d.\,h. für jeden positiven ganzzahligen Wert von $\mu$, $\varphi'(\mu-1)N=(1-\mu)\varphi'(-N)$, woraus sich leicht ergibt, dass $\frac{\varphi'\Delta}{\Delta}$ im Allgemeinen eine konstante Größe sein muss, die wir mit $k$ bezeichnen. Daraus folgt $\log\varphi\Delta=\tfrac12 k\Delta\Delta+\text{Konst.}$ oder, wenn wir die Basis der hyperbolischen Logarithmen mit $e$ bezeichnen und $\text{Konst.}=\log\varkappa$ annehmen,
\[
\varphi\Delta=\varkappa e^{\frac12 k\Delta\Delta}.
\]
Ferner ist leicht zu sehen, dass $k$ notwendig negativ sein muss, damit $\Omega$ tatsächlich ein Maximum annehmen kann, weshalb wir $\tfrac12 k=-hh$ setzen; und da nach dem eleganten Satz, den zuerst der berühmte \textsc{Laplace} gefunden hat, das Integral $\int e^{-hh\Delta\Delta}\mathrm{d}\Delta$ von $\Delta=-\infty$ bis $\Delta=+\infty$ den Wert $=\frac{\BareRoot{\pi}}{h}\mathpunct{\raisebox{.45ex}{,}}$ annimmt (wobei $\pi$ den halben Umfang eines Kreises mit dem Radius $1$ bezeichnet), wird unsere Funktion zu
\[
\varphi\Delta=\frac{h}{\BareRoot{\pi}}e^{-hh\Delta\Delta}.
\]
